\section*{Bab \ref{ch:g8:sound-pitch-loudness} --- Bunyi: Kelajuan, Nada, Kenyaringan}

\begin{solution}{exo:g8:sound-pitch-loudness:1}
\omterm{def:g8:sound-pitch-loudness:frequency}{Frekuensi}: banyaknya ayunan lengkap tiap sekon, dalam \omterm{def:g8:sound-pitch-loudness:frequency}{hertz}
(\unit{Hz}). Nada A konser: \qty{440}{Hz}; oktaf di atasnya:
\qty{880}{Hz} --- digandakan.
\end{solution}

\begin{solution}{exo:g8:sound-pitch-loudness:2}
$4$ ayunan dalam \qty{0.02}{s}: \qty{200}{Hz}; $16$ ayunan:
\qty{800}{Hz} --- yang kedua terdengar lebih tinggi (dua oktaf,
sebenarnya).
\end{solution}

\begin{solution}{exo:g8:sound-pitch-loudness:3}
Kira-kira \qty{20}{Hz} sampai \qty{20000}{Hz}. Di bawahnya:
\omterm{prop:g8:sound-pitch-loudness:range}{infrasonik} --- gemuruh jarak jauh milik gajah. Di atasnya: \omterm{prop:g8:sound-pitch-loudness:range}{ultrasonik}
--- detak berburu milik kelelawar.
\end{solution}

\begin{solution}{exo:g8:sound-pitch-loudness:4}
Lebih nyaring, nada yang sama: puncaknya bertumbuh lebih tinggi,
jaraknya tidak berubah. Lebih tinggi, kenyaringan yang sama: puncaknya
berdesak lebih rapat, tingginya tidak berubah. Dua tombol yang mandiri.
\end{solution}

\begin{solution}{exo:g8:sound-pitch-loudness:5}
Masing-masing menaikkan \omterm{def:g8:sound-pitch-loudness:frequency}{frekuensinya}: dawai yang lebih pendek, lebih
tegang, dan lebih tipis semuanya bergetar lebih banyak kali tiap sekon
--- $f$ yang lebih tinggi, nada yang lebih tinggi.
\end{solution}

\begin{solution}{exo:g8:sound-pitch-loudness:6}
Daun \qty{10}{dB}; lalu lintas \qty{80}{dB}; barisan depan
\qty{110}{dB}; rasa sakit dekat \qty{120}{dB}. Pita kerusakannya
terbuka sekitar \qty{85}{dB}.
\end{solution}

\begin{solution}{exo:g8:sound-pitch-loudness:7}
Nada siulannya duduk tepat di atas \qty{20000}{Hz} --- di luar jendela
pemiliknya, di dalam jendela anjingnya. Rusak bagi sepasang telinga,
menusuk bagi yang lain: jendelanya berbeda menurut jenis makhluknya.
\end{solution}

\begin{solution}{exo:g8:sound-pitch-loudness:8}
Bahwa \omterm{def:g7:motion-average-speed:formula}{kelajuannya} satu dan sama bagi setiap nada dan setiap
kenyaringan: \omterm{def:g7:motion-average-speed:formula}{kelajuan} yang berbeda akan mengaburkan kunci nada yang
jauh menjadi arpegio lalu memperlambat yang lirih di belakang yang
nyaring --- konser tamannya membantah keduanya tiap malam.
\end{solution}

\begin{solution}{exo:g8:sound-pitch-loudness:9}
$50000 \times 0.01 = 500$ ayunan tiap detaknya. Lukisan yang halus
memerlukan kuas yang halus: hanya \omterm{def:g3:sound-around-us:vibration}{getaran} dengan ayunan yang sangat
pendek --- $f$ yang sangat tinggi --- yang dapat meraba perincian yang
kecil di dalam \omterm{ex:g4:how-sound-travels:echo}{gemanya}; gemuruh berfrekuensi rendah yang kasar
menyapu lewat seekor ngengat tanpa memperhatikannya.
\end{solution}

\begin{solution}{exo:g8:sound-pitch-loudness:10}
Kedua tutsnya duduk di dalam jendelanya --- \qty{27}{Hz} tepat di atas
lantai \qty{20}{Hz}-nya, \qty{4200}{Hz} dengan nyaman di bawah
langit-langitnya. Tetapi usia mengikis jendelanya dari \emph{atas}:
telinga penyetelnya sendiri menumpul lebih dahulu pada nada tingginya,
sehingga penganalisisnya dipanggil untuk oktaf yang tertinggi sebelum
yang terendah.
\end{solution}

\begin{solution}{exo:g8:sound-pitch-loudness:11}
Pertumbuhan sepuluh kali lipat yang bersembunyi di dalam langkah yang
sama --- tiap beberapa anak tangga mengalikan kekuatan \omterm{def:g3:sound-around-us:sound}{bunyinya}, bukan
menambahinya. Jadi ``sedikit lebih nyaring'' pada tangganya berarti
\emph{jauh} lebih kuat bagi telinganya: anggarannya runtuh dari
berjam-jam menjadi beberapa menit sepanjang apa yang terbaca sebagai
putaran tombol yang kecil.
\end{solution}

\begin{solution}{exo:g8:sound-pitch-loudness:12}
Pemeriksaan pianonya: penganalisisnya di samping dawainya, lalu
bandingkan petunjuknya dengan $440$. Kenyaringan barisan belakangnya:
berdirilah di sana dengan petunjuk \omterm{ex:g8:sound-pitch-loudness:decibels}{desibelnya} selama uji gendangnya.
Peranginannya: bidikkan tampilan \omterm{prop:g8:colors-spectra:mixture}{spektrum} penganalisisnya ke ujung
yang rendah lalu carilah puncak di bawah \qty{20}{Hz}. Yang paling
tidak dapat dipercaya: putusan \omterm{prop:g8:sound-pitch-loudness:range}{infrasoniknya} --- mikrofon telepon
dibangun untuk suara orang lalu menjadi tuli dekat lantai jendelanya,
sehingga layar yang sunyi tidak membuktikan banyak di bawah sana.
\end{solution}

\begin{solution}{pb:g8:sound-pitch-loudness:1}
\textbf{1.} Pianonya terlalu rendah --- setiap nada A bergetar delapan
ayunan tiap sekon terlalu sedikit. Tombol nadanya: kunci seorang
penyetel, bukan sebuah penguat.
\textbf{2.} Angka \qty{41}{Hz} duduk tepat di atas lantai
\qty{20}{Hz} jendelanya: terdengar, dengan dalam --- dan dengan ayunan
yang selambat itu, gendang dadanya sendiri ikut menjawab: terasa
sebanyak terdengar.
\textbf{3.} Nada yang sama (jarak yang sama, \qty{440}{Hz}), dengan
\omterm{prop:g8:sound-pitch-loudness:loudness}{amplitudo} tiga kali lipat: penyanyinya menyanyikan nada garpu talanya,
jauh lebih nyaring.
\textbf{4.} Hanya telinga yang termuda --- langit-langit jendelanya
dekat \qty{20000}{Hz} menjadi milik anak-anak; jendela kebanyakan
orang dewasa sudah menurun, dan tidak ada pendengar yang mengesahkan
\qty{22000}{Hz}.
\textbf{5.} Bagian roknya (\qty{100}{dB}) dan babak penutupnya
(\qty{110}{dB}) keduanya berdiri di atas garis kerusakan
\qty{85}{dB}: wilayah beberapa menit saja; bagian akustiknya lulus.
\textbf{6.} Kira-kira \qty{80}{dB} di balik penyumbatnya --- kembali
di bawah garis kerusakannya: paparan sesore dibuat dapat dilalui, dan
itulah sebabnya yang bijak selalu memakainya.
\textbf{7.} Kulit yang menyebar itu: tiap perasannya dibagikan atas
kulit yang kian luas menurut jarak --- kenyaringannya (\omterm{prop:g8:sound-pitch-loudness:loudness}{amplitudonya})
memudar. Nadanya menunggangi tanpa berubah: jarak memutar satu tombol
saja.
\textbf{8.} Bermain dengan lembut memutar tombol \omterm{prop:g8:sound-pitch-loudness:loudness}{amplitudonya} saja ---
\omterm{def:g8:sound-pitch-loudness:frequency}{frekuensi} setiap nadanya, tinggi dan rendah, tinggal persis di
tempatnya. (``Lebih tumpul''-nya adalah muslihat telinganya sendiri
pada kenyaringan yang rendah, bukan fisika dawainya.)
\textbf{9.} $680 \div 340 = \qty{2}{s}$ dari kilat ke dentum. Dua
gedebuk tiap sekon: \qty{2}{Hz} --- jauh di bawah lantainya: terasa di
dada, tidak pernah terdengar sebagai nada.
\textbf{10.} Ayunan terendah basnya: dekat lantai jendelanya dan di
bawahnya, tepat di tempat mikrofon yang dibangun untuk suara orang
menjadi tuli --- gemuruhnya nyata; teleponnya tidak pernah
menangkapnya.
\textbf{11.} Misalnya: ``Nada adalah \omterm{def:g8:sound-pitch-loudness:frequency}{frekuensi} --- hitungan ayunan
tiap sekon, dan seluruh urusan sang penyetel. Kenyaringan adalah
\omterm{prop:g8:sound-pitch-loudness:loudness}{amplitudo} --- ukuran ayunannya, seluruh buku sang perawat. Dan tidak
satu pun tombolnya menyentuh \omterm{def:g7:motion-average-speed:formula}{kelajuannya}: setiap nada malam ini
menyeberangi aulanya pada \qty{340}{m/s} yang sama.''
\textbf{12.} Proposisi jendelanya: pengusirnya bernada di atas
\qty{20000}{Hz} --- di luar setiap jendela manusia, di dalam jendela
tikusnya. Satu \omterm{def:g3:sound-around-us:vibration}{getaran}, dua khalayak, tanpa pertentangan.
\end{solution}
